% answers_short.tex —— 2013 年·上海卷·文科 · 简短答案速查
% 由 tools/make_short_answers.py 自 answers.tex 的 \gksolhead/\gkanswer 提取，
% 并与 2013/2013_数学真题_上海文_答案要点.md 逐题比对后生成，勿手改（改源 MD 后重跑）。

\begin{gknotice}
\textbf{使用说明}：本节只给最终答案便于快速核对，完整推导见解析版。答案为编写者独立推导，\textbf{非官方标准答案}。
\end{gknotice}

\gkpart{填空题}
\begin{itemize}[leftmargin=2.6em]
  \item[\textbf{1.}] $\left(0,\dfrac12\right)$。由 $\dfrac{x}{2x-1}<0$ 得 $x(2x-1)<0$，两根为 $0$ 与 $\dfrac12$，故 $0<x<\dfrac12$。
  \item[\textbf{2.}] $15$。由 $a_1+a_4=a_2+a_3$ 得 $2(a_2+a_3)=30$，故 $a_2+a_3=15$。
  \item[\textbf{3.}] $-2$。纯虚数要求实部为 $0$ 且虚部不为 $0$：$m^2+m-2=0$ 给出 $m=1$ 或 $m=-2$；$m^2-1\neq0$ 排除 $m=\pm1$，故 $m=-2$。
  \item[\textbf{4.}] $3$。由 $x-2=0$ 得 $x=2$；由 $x-y=1$ 得 $y=1$，故 $x+y=3$。
  \item[\textbf{5.}] $\dfrac{2\pi}3$。由 $c^2=a^2+b^2+ab$ 与余弦定理 $c^2=a^2+b^2-2ab\cos C$ 得 $\cos C=-\dfrac12$，故 $C=\dfrac{2\pi}3$。
  \item[\textbf{6.}] $78$。按人数加权：$0.4\times75+0.6\times80=30+48=78$。
  \item[\textbf{7.}] $-2$。通项为 $\mathrm C_5^k a^k x^{10-3k}$，令 $10-3k=7$ 得 $k=1$，故 $5a=-10$，$a=-2$。
  \item[\textbf{8.}] $\log_3 4$。令 $t=3^x>0$（$t\neq1$），方程化为 $\dfrac{9}{t-1}=t-1$，即 $(t-1)^2=9$，$t=4$（$t=-2$ 舍去），故 $x=\log_3 4$。
  \item[\textbf{9.}] $-\dfrac79$。由已知 $\cos(x-y)=\dfrac13$，故 $\cos(2x-2y)=2\cos^2(x-y)-1=\dfrac29-1=-\dfrac79$。
  \item[\textbf{10.}] $\sqrt3$。取 $O(0,0,l)$，$A(r,0,0)$，则 $\overrightarrow{OA}=(r,0,-l)$，而 $\overrightarrow{BC}=(0,0,l)$，故 $\cos\dfrac\pi6=\dfrac{l}{\sqrt{r^2+l^2}}=\dfrac{\sqrt3}2$，得 $l^2=3r^2$，$\dfrac lr=\sqrt3$。
  \item[\textbf{11.}] $\dfrac57$。总取法 $\mathrm C_7^2=21$；两球编号之积为奇数须同为奇数，共 $\mathrm C_4^2=6$ 种，故积为偶数的概率为 $1-\dfrac6{21}=\dfrac57$。
  \item[\textbf{12.}] $\dfrac{4\sqrt6}3$。设椭圆为 $\dfrac{x^2}4+\dfrac{y^2}{b^2}=1$，$B(2,0)$，则 $\angle CBA=\dfrac\pi4$ 且 $BC=\sqrt2$ 给出 $C(1,\pm1)$；代入得 $\dfrac14+\dfrac1{b^2}=1$，$b^2=\dfrac43$，于是 $c^2=4-\dfrac43=\dfrac83$，两焦点距离 $2c=\dfrac{4\sqrt6}3$。
  \item[\textbf{13.}] $\left[\dfrac15,+\infty\right)$。由均值不等式 $9x+\dfrac{a^2}x\geqslant2\sqrt{9a^2}=6a$（等号在 $x=\dfrac a3$ 取得），故只需 $6a\geqslant a+1$，即 $a\geqslant\dfrac15$。
  \item[\textbf{14.}] $-5$。
\end{itemize}

\gkpart{选择题}
\begin{center}\begin{tabular}{ccccc}
\toprule
题号 & 15 & 16 & 17 & 18 \\
\midrule
答案 & A & B & B & D \\
\bottomrule
\end{tabular}\end{center}

\gkpart{解答题}
\begin{itemize}[leftmargin=2.6em]
  \item[\textbf{19.}] $V=\dfrac{\sqrt3}3$，$S=3\sqrt3$。
  \item[\textbf{20.}] (1) 证明见解析；(2) $x=6$ 千克/小时，最大利润 $457500$ 元
  \item[\textbf{21.}] (1) 既不是奇函数也不是偶函数；(2) $20$ 或 $21$
  \item[\textbf{22.}] (1) $a_2=2$，$a_3=0$，$a_4=2$；(2) $a_1=1$ 或 $a_1=2+\sqrt2$；(3) 存在，$a_1=1$
  \item[\textbf{23.}] (1) $x=-\sqrt3$（答案不唯一）；(2) 证明见解析；(3) 证明见解析
\end{itemize}
